% ============================================================================
\section{微分形式与模形式}
\label{sec:differentials-and-modforms}
\label{sec:modular-forms-differentials}
% ============================================================================

\textbf{为什么模形式会和微分形式扯上关系？}

到目前为止，模形式看起来像是一类满足特殊变换律的函数。
但这种定义虽然能算，几何味道却还不够强。
如果我们想真正理解模形式为什么和亏格、除子、Riemann--Roch 连在一起，
最自然的办法就是把它们改写成黎曼面上的几何对象。

\textbf{最关键的观察发生在权 $2$。}
一个形如
\[
  \omega=f(\tau)\,d\tau
\]
的 $1$-形式在变量替换下会多出一个 Jacobian 因子。
而这个 Jacobian 因子恰好是 $(c\tau+d)^{-2}$；
所以要让 $\omega$ 对 $\Gamma$ 不变，
$f$ 就必须满足权 $2$ 的模变换律。
换句话说：\emph{权 $2$ 根本不是偶然冒出来的数字，它正是微分 $d\tau$ 的几何权重。}

\textbf{这一点一旦看懂，就立刻得到一个极其漂亮的结论。}
权 $2$ cusp form 与 $X(\Gamma)$ 上的全纯 $1$-形式完全等价。
因此
\[
  \dim \Sk_2(\Gamma)=g\bigl(X(\Gamma)\bigr).
\]
这就是解析空间维数和曲面亏格之间最直接、最优雅的桥梁。

\begin{definition}[上半平面上的 $1$-形式]
\label{def:one-form-on-h}
在 $\HH$ 上，一个全纯 $1$-形式写作
\[
  \omega=f(\tau)\,d\tau,
\]
其中 $f$ 是 $\HH$ 上的全纯函数。
若 $\gamma=\begin{pmatrix}a&b\\c&d\end{pmatrix}\in\SL_2(\R)$，则
\[
  \gamma\tau=\frac{a\tau+b}{c\tau+d},
  \qquad
  d(\gamma\tau)=\frac{d\tau}{(c\tau+d)^2}.
\]
\end{definition}

\begin{proposition}[下降到商空间的条件]
\label{prop:descent-weight-2}
设 $\omega=f(\tau)\,d\tau$ 是 $\HH$ 上的全纯 $1$-形式。
则 $\omega$ 对 $\Gamma$ 不变，也就是
\[
  \gamma^*\omega=\omega \qquad (\gamma\in\Gamma),
\]
当且仅当
\[
  f(\gamma\tau)=(c\tau+d)^2f(\tau)
\]
对所有 $\gamma=\begin{pmatrix}a&b\\c&d\end{pmatrix}\in\Gamma$ 成立。
也就是说，$f$ 恰好是权 $2$ 的模形式。
\end{proposition}

\begin{proof}
由拉回公式，
\[
  \gamma^*\omega=f(\gamma\tau)\,d(\gamma\tau)
  =f(\gamma\tau)\frac{d\tau}{(c\tau+d)^2}.
\]
要让它等于 $f(\tau)\,d\tau$，
恰好等价于
\[
  f(\gamma\tau)=(c\tau+d)^2f(\tau).
\qedhere
\]
\end{proof}

\begin{theorem}
\label{thm:s2-holomorphic-differentials}
设 $\Gamma$ 为有限指数子群并且 $-I\in\Gamma$。
则赋值
\[
  f \longmapsto \omega_f=f(\tau)\,d\tau
\]
给出线性同构
\[
  \Sk_2(\Gamma)\cong H^0\bigl(X(\Gamma),\Omega^1\bigr),
\]
其中右边是紧模曲线 $X(\Gamma)$ 上全纯 $1$-形式的空间。
\end{theorem}

\begin{proof}
由\cref{prop:descent-weight-2}，权 $2$ 条件保证 $f(\tau)d\tau$ 能从 $\HH$ 下降到
$Y(\Gamma)$ 上。
还差最后一步：说明它在每个尖点处也能全纯延拓到紧化后的 $X(\Gamma)$。

先看尖点 $\infty$，设其宽度为 $h$，局部坐标为
\[
  q_h=e^{2\pi i\tau/h}.
\]
则
\[
  d\tau=\frac{h}{2\pi i}\frac{dq_h}{q_h}.
\]
于是
\[
  \omega_f=f(\tau)\,d\tau
  =\frac{h}{2\pi i}f(\tau(q_h))\frac{dq_h}{q_h}.
\]
若 $f$ 是 cusp form，则其 $q_h$-展开没有常数项，即
$f(\tau(q_h))=a_1q_h+a_2q_h^2+\cdots$。
因此上式变成
\[
  \omega_f=\frac{h}{2\pi i}(a_1+a_2q_h+\cdots)dq_h,
\]
在 $q_h=0$ 处全纯。
对一般尖点 $s=\sigma\infty$，把 $f$ 换成 $f|_2\sigma$ 完全同理。
所以 $f\mapsto \omega_f$ 的像确实落在
$H^0\bigl(X(\Gamma),\Omega^1\bigr)$ 中。

反过来，若 $\omega$ 是 $X(\Gamma)$ 上的全纯 $1$-形式，
把它拉回到 $\HH$，便可写成 $f(\tau)d\tau$。
拉回后的 $\Gamma$-不变性迫使 $f$ 满足权 $2$ 变换律；
而 $\omega$ 在尖点处全纯，又迫使 $f$ 在每个尖点的常数项为 $0$。
故 $f\in\Sk_2(\Gamma)$。
线性性显然，于是得到同构。
\end{proof}

\begin{corollary}
\label{cor:dim-s2-genus}
\[
  \dim \Sk_2(\Gamma)=g\bigl(X(\Gamma)\bigr).
\]
\end{corollary}

\begin{proof}
紧黎曼面上全纯 $1$-形式空间的维数等于其亏格。
结合\cref{thm:s2-holomorphic-differentials} 即得。
\end{proof}

\begin{example}[$X_0(11)$ 上唯一的权 $2$ cusp form]
\label{ex:x011-weight2}
由\cref{ex:genus-small-levels,cor:dim-s2-genus}，
$g\bigl(X_0(11)\bigr)=1$，故
\[
  \dim \Sk_2\bigl(\Gamma_0(11)\bigr)=1.
\]
因此 $\Sk_2\bigl(\Gamma_0(11)\bigr)$ 由一条非零形式张成；
它的归一化生成元可以写成
\[
  f(\tau)=q-2q^2-q^3+2q^4+q^5+\cdots.
\]
对应的微分形式 $f(\tau)d\tau$ 就是 $X_0(11)$ 上的全纯 $1$-形式，
这也解释了为什么 $X_0(11)$ 本身是一条椭圆曲线。
\end{example}

对于更高偶权 $k$，同样的变换计算告诉我们：
若 $f\in\Mk_k(\Gamma)$，则
\[
  f(\tau)(d\tau)^{k/2}
\]
会下降成 $X(\Gamma)$ 上规范丛 $\Omega$ 的 $k/2$ 次张量幂
$\Omega^{\otimes k/2}$ 的一个局部截面。
若 $f$ 在各尖点也全纯，则它就是一个全纯截面；
若还在尖点消失，则对应的零阶数更高。

% figures/ch02-modforms-differentials.tex
% 模形式与微分形式 / 线丛截面的对应
\begin{figure}[ht]
\centering
\begin{tikzpicture}[
  >=Stealth,
  font=\small,
  box/.style={draw, rounded corners=4pt, minimum width=3.5cm, minimum height=1.15cm, align=center},
  arrow/.style={->, thick}
]
  \node[box, fill=blue!10] (mk) at (0,2.2)
    {$f\in \Mk_k(\Gamma)$\\满足 $f(\gamma\tau)=(c\tau+d)^k f(\tau)$};
  \node[box, fill=green!10] (tensor) at (5.3,2.2)
    {$f(\tau)(d\tau)^{k/2}$\\看作 $\Omega^{\otimes k/2}$ 的局部截面};
  \node[box, fill=orange!12] (forms) at (5.3,0)
    {$H^0\bigl(X(\Gamma),\Omega^{\otimes k/2}\bigr)$\\全纯截面空间};
  \node[box, fill=red!10] (k2) at (0,0)
    {$f\in \Sk_2(\Gamma)$\\在每个尖点都消失};

  \draw[arrow] (mk) -- node[above] {几何化} (tensor);
  \draw[arrow] (tensor) -- node[right] {下降到 $X(\Gamma)$} (forms);
  \draw[arrow] (k2) -- node[below] {$\omega=f(\tau)d\tau$} (forms);

  \node[align=center] at (2.65,-1.65)
    {特别地，当 $k=2$ 时，尖点形式就是全纯 $1$-形式；\\
     因而 $\dim \Sk_2(\Gamma)=g\bigl(X(\Gamma)\bigr)$.};
\end{tikzpicture}
\caption{模形式的几何解释。权 $2$ 时得到真正的微分形式；更高偶权时得到规范丛张量幂的全纯截面。}
\label{fig:modforms-differentials}
\end{figure}


因此，Riemann--Roch 可以把模形式空间的维数写成几何维数公式。
对含 $-I$ 的群，一个典型的前瞻性公式是
\[
  \dim \Mk_k(\Gamma)
  =(k-1)(g-1)+\left\lfloor\frac{k}{4}\right\rfloor\nu_2
   +\left\lfloor\frac{k}{3}\right\rfloor\nu_3
   +\left(\frac{k}{2}-1\right)\nu_\infty,
\]
再配上常数项修正，就会得到\cref{ch:dimension} 中真正的维数公式。
这里我们先记住它传达的核心信息：
\emph{模形式的维数，本质上是一条紧黎曼面上线丛截面空间的维数。}
